\clearpage\thispagestyle{empty}
\pdfbookmark[0]{Source for the 6 August 1808 Letter to von Lindenau}{lindenau1808aug06-sources}
\begin{center}\textbf{Source for the 6 August 1808 Letter to von Lindenau}\end{center}
\noindent This AI translation follows the 6 August 1808 letter from Carl Friedrich Gauss to Bernhard August von Lindenau in Gauss’s \href{https://gdz.sub.uni-goettingen.de/id/PPN236007467}{\textit{Werke}, vol.~VI} (Göttingen, 1874), pp.~305--306. The first printing appeared as item~XXI, \textit{Auszug aus einem Schreiben des Hrn. Dr. Gauß}, in the August issue of \textit{Monatliche Correspondenz}, vol.~XVIII, pp.~\href{https://www.digitale-sammlungen.de/en/view/bsb10538609?page=193}{173}, \href{https://www.digitale-sammlungen.de/en/view/bsb10538609?page=194}{174} and \href{https://www.digitale-sammlungen.de/en/view/bsb10538609?page=195}{175}.
\par\medskip
The collected edition includes the handwritten Ceres entry and refers readers to the November notice for Vesta’s observations. Ernst Schering cites the related notice in the Göttingische gelehrte Anzeigen, issue 127, \href{https://manifests.sub.uni-goettingen.de/iiif/image/gdz:PPN319721507_1808:00001271/full/1600,/0/default.jpg}{1265}, for the sentence about the instruments still in Paris. Gauss’s Latin phrase \textit{levi calamo} conveys that he made the calculation quickly and without ceremony. For the instrument terminology, see \href{https://www.gutenberg.org/files/53172/old/53172-h/53172-h.htm#fig163}{J. Norman Lockyer, \textit{Stargazing: Past and Present}, p.~368}. Translation and typesetting produced by AI.
\clearpage
